\section{API de resultados}

\subsection{Descripcion}

Este documento define la integracion conceptual con la API externa de resultados. Su objetivo es dejar claro que espera el sistema de un proveedor externo, como se adapta esa informacion al dominio interno y que limites debe respetar la integracion para no contaminar reglas del juego con el modelo del proveedor.

\subsection{Alcance}

Aplica a la integracion de resultados de partidos y metadatos asociados necesarios para soportar cierre funcional del partido y trazabilidad. Para \texttt{QH-111} estas decisiones quedan cerradas a nivel de proveedor, endpoints validados, mapping, normalizacion y manejo de errores. El objetivo es dejar \texttt{QH-111} decision-complete antes de implementacion.

\subsection{Definiciones}

\begin{itemize}
  \item \textbf{Proveedor de resultados}: servicio externo que entrega datos oficiales o candidatos a oficiales sobre partidos.
  \item \textbf{Adaptacion}: transformacion del dato externo al modelo interno del dominio.
  \item \textbf{Resultado vigente}: valor oficial que el sistema reconoce para scoring y vistas derivadas.
\end{itemize}

\subsection{Reglas}

\begin{itemize}
  \item El modelo interno del dominio es la referencia del sistema; la API externa debe adaptarse a ese modelo y no al reves.
  \item Todo dato recibido desde el proveedor debe poder rastrearse hasta su fuente, momento de recepcion y partido interno asociado.
  \item La API externa no se considera autoridad absoluta; su salida puede ser validada, ignorada o reemplazada por operacion manual auditada.
  \item El sistema solo debe consumir del proveedor la informacion necesaria para operar la quiniela y no depender de campos irrelevantes del modelo externo.
\end{itemize}

\subsubsection{Datos minimos esperados del proveedor}

\begin{itemize}
  \item identificador externo del partido;
  \item fecha y hora oficial del kickoff;
  \item equipos involucrados;
  \item estado funcional del partido o equivalente interpretable;
  \item goles oficiales utilizable para la regla de 90 minutos mas agregado;
  \item marca temporal de actualizacion del dato externo.
\end{itemize}

\subsection{Proveedor cerrado para MVP (QH-111)}

\subsubsection{Proveedor recomendado}

Proveedor principal para MVP Mundial 2026: \textbf{football-data.org}.

Motivo de decision (basado en POC, no en supuestos):
\begin{itemize}
  \item football-data.org permitio consultar directamente World Cup 2026 en el plan probado.
  \item \texttt{GET /competitions/WC/matches?season=2026} devolvio los 104 partidos del Mundial 2026.
  \item Permite match individual \texttt{GET /matches/\{matchId\}} y batch por ids \texttt{GET /matches?ids=...}.
  \item API-Football (API-Sports) quedo como alternativa tecnicamente viable, pero en plan Free bloqueo \texttt{league=1\&season=2026} y batch IDs; quedo condicionada a plan Pro.
\end{itemize}

\subsubsection{Base URL, auth y variables de entorno}

Base URL:
\begin{itemize}
  \item \texttt{https://api.football-data.org/v4}
\end{itemize}

Auth:
\begin{itemize}
  \item header \texttt{X-Auth-Token: <FOOTBALL\_DATA\_API\_TOKEN>}
\end{itemize}

Variables recomendadas:
\begin{itemize}
  \item \texttt{FOOTBALL\_DATA\_BASE\_URL=https://api.football-data.org/v4}
  \item \texttt{FOOTBALL\_DATA\_API\_TOKEN=<secret>}
  \item \texttt{FOOTBALL\_DATA\_COMPETITION\_CODE=WC}
  \item \texttt{FOOTBALL\_DATA\_TARGET\_SEASON=2026}
  \item \texttt{FOOTBALL\_DATA\_PROVIDER\_NAME=football-data.org}
\end{itemize}

\subsubsection{Endpoints validados (POC)}

Los siguientes endpoints fueron validados durante POC con token valido:
\begin{itemize}
  \item \texttt{GET /competitions} (token valido, WC disponible)
  \item \texttt{GET /competitions/WC} (currentSeason=2026, fechas del torneo)
  \item \texttt{GET /competitions/WC/matches?season=2026} (104 matches, campos utiles)
  \item \texttt{GET /matches?dateFrom=YYYY-MM-DD\&dateTo=YYYY-MM-DD} (fallback por ventana)
  \item \texttt{GET /matches?dateFrom=...\&dateTo=...\&status=SCHEDULED} (filtro para futuros; response puede traer \texttt{status=TIMED})
  \item \texttt{GET /matches/\{matchId\}} (match individual)
  \item \texttt{GET /matches?ids=id1,id2,id3} (batch por ids)
  \item \texttt{GET /competitions/WC/teams?season=2026} (48 equipos)
\end{itemize}

No aplica para World Cup como CUP:
\begin{itemize}
  \item \texttt{GET /competitions/WC/standings?season=2026} (404 / no aplicable)
\end{itemize}

Nota importante de POC:
\begin{itemize}
  \item En eliminatoria, football-data.org puede devolver \texttt{homeTeam} o \texttt{awayTeam} como \texttt{null} en partidos futuros. La implementacion no debe asumir equipos definidos desde el inicio.
\end{itemize}

\subsubsection{Errores validados (POC)}

\begin{itemize}
  \item Match inexistente: \texttt{GET /matches/999999999} devuelve error esperado con \texttt{errorCode=400}. Debe mapearse a \texttt{external\_match\_not\_found}.
  \item Token invalido: \texttt{GET /competitions} con token invalido devuelve error esperado con \texttt{errorCode=400}. Debe mapearse a \texttt{provider\_auth\_invalid} (no transitorio).
\end{itemize}

\subsection{Contrato funcional cerrado para QH-111}

\subsubsection{Encapsulamiento (adapter)}

\begin{itemize}
  \item La integracion debe estar encapsulada detras de interfaz interna (provider/adapter).
  \item Implementacion recomendada: \texttt{FootballDataResultsProviderAdapter}.
  \item El caso de uso no debe depender del payload crudo del proveedor.
\end{itemize}

\subsubsection{DTO normalizado minimo (salida del adapter)}

El adapter debe producir un DTO interno normalizado con al menos:
\begin{itemize}
  \item \texttt{provider}
  \item \texttt{externalMatchId}
  \item \texttt{providerStatus}
  \item \texttt{normalizedStatus}
  \item \texttt{homeGoals}
  \item \texttt{awayGoals}
  \item \texttt{utcDate}
  \item \texttt{stage}
  \item \texttt{group}
  \item \texttt{matchday}
  \item \texttt{homeTeamExternalId}
  \item \texttt{awayTeamExternalId}
  \item \texttt{providerUpdatedAt} (o equivalente si esta disponible)
  \item \texttt{externalReference}
\end{itemize}

\subsubsection{Adaptacion al dominio interno}

\begin{itemize}
  \item El proveedor puede usar nomenclaturas de estado distintas; el backend debe traducirlas a los estados funcionales internos del proyecto.
  \item El sistema debe identificar un partido interno de forma unica antes de aceptar un resultado externo como candidato valido.
  \item Si el proveedor incluye informacion no utilizable por el reglamento, como tiempo extra o penales, esa informacion no debe contaminar el resultado oficial usado para scoring.
  \item La adaptacion debe producir un valor interpretable por el dominio aun cuando el proveedor entregue estructuras mas complejas que las necesarias para la quiniela.
\end{itemize}

\subsubsection{Criterios de aceptacion de un resultado externo}

\begin{itemize}
  \item el partido interno asociado existe;
  \item el dato puede mapearse al estado funcional del partido;
  \item los goles y el contexto permiten determinar un resultado compatible con el reglamento;
  \item no existe override manual vigente que impida adoptar automaticamente el valor de API.
\end{itemize}

\subsubsection{Condicion para persistir resultado oficial}

\begin{itemize}
  \item Solo persistir cuando \texttt{providerStatus = FINISHED}.
  \item Adicionalmente, \texttt{score.fullTime.home} y \texttt{score.fullTime.away} deben ser no \texttt{null}.
  \item Fuente del marcador oficial: \texttt{score.fullTime.home} y \texttt{score.fullTime.away}.
  \item No persistir resultados desde estados intermedios o futuros.
\end{itemize}

\subsubsection{Traduccion de estados externos (football-data.org)}

Mapeo minimo:
\begin{itemize}
  \item \texttt{TIMED} $\rightarrow$ programado / no persistible
  \item \texttt{SCHEDULED} (como query param externo) $\rightarrow$ programado / no persistible
  \item \texttt{FINISHED} $\rightarrow$ persistible si \texttt{score.fullTime} no es \texttt{null}
  \item \texttt{LIVE}, \texttt{IN\_PLAY}, \texttt{PAUSED} $\rightarrow$ skip / no persistible (MVP sin live sync)
  \item \texttt{POSTPONED}, \texttt{SUSPENDED}, \texttt{CANCELLED} $\rightarrow$ skip / no persistible
  \item estado desconocido $\rightarrow$ skip seguro + auditoria
\end{itemize}

\subsubsection{Mapping externo/interno (schema recomendado)}

\textbf{Match:} mapping directo por proveedor + id externo. No usar \texttt{Result.externalReference} para localizar partidos.

Recomendacion de schema para \texttt{Match}:
\begin{verbatim}
externalProvider String?
externalMatchId  String?

@@unique([externalProvider, externalMatchId])
@@index([externalProvider])
\end{verbatim}

Politica ante partido no mapeado:
\begin{itemize}
  \item Si llega un resultado externo cuyo \texttt{externalMatchId} no corresponde a un \texttt{Match} interno sincronizable, hacer skip seguro.
  \item No crear \texttt{Match} automaticamente en QH-111.
  \item Registrar evidencia con razon \texttt{external\_match\_not\_mapped}.
\end{itemize}

\textbf{Team:} mapping auxiliar por \texttt{team.id} del proveedor.

Recomendacion de schema para \texttt{Team}:
\begin{verbatim}
externalProvider String?
externalTeamId   String?
shortName        String?
crestUrl         String?

@@unique([externalProvider, externalTeamId])
\end{verbatim}

\subsection{Limites explicitos de QH-111}

\begin{itemize}
  \item No incluye cron programado (solo caso de uso interno + superficie manual/admin ``Sync now'').
  \item No incluye recomputacion real de puntos.
  \item No incluye creacion automatica de \texttt{Match} desde proveedor.
  \item No incluye creacion automatica de \texttt{Team} desde proveedor (solo mapping auxiliar si existe ticket separado que lo autorice).
  \item No incluye soporte multi-provider activo.
  \item No incluye tabla \texttt{MatchExternalMapping} ni \texttt{TeamExternalMapping}.
  \item No incluye \texttt{ResultHistory}.
  \item No incluye live sync / polling near-real-time.
  \item No incluye exposicion publica de metadata externa.
\end{itemize}

\subsection{Efecto en lectura publica}

La lectura publica de resultado por partido (ej. \texttt{GET /matches/:matchId/result}) debe reflejar inmediatamente el \texttt{Result} persistido.

La respuesta publica no debe exponer:
\begin{itemize}
  \item \texttt{source}
  \item \texttt{isOfficial}
  \item \texttt{externalReference}
  \item \texttt{recordedByUserId}
  \item payload crudo del proveedor
\end{itemize}

\subsubsection{Superficies internas consumidoras}

\begin{itemize}
  \item modulo \texttt{results}: adapta y conserva referencia de fuente;
  \item modulo \texttt{matches}: actualiza estado funcional del partido cuando aplique;
  \item modulo \texttt{scoring}: consume solo el resultado vigente ya adaptado;
  \item frontend: consume el resultado visible oficial y nunca el payload externo crudo.
\end{itemize}

\subsection{Edge Cases}

\begin{itemize}
  \item Si el proveedor entrega un resultado parcial o transitorio, ese dato no debe consolidarse como resultado final sin pasar por la logica interna de aceptacion.
  \item Si el proveedor cambia el kickoff oficial, el sistema debe tratar ese cambio como dato integrable del partido y no como simple texto informativo.
  \item Si el proveedor no puede mapear limpiamente el partido o su estado, el dato debe quedar rechazado o pendiente de revision operativa.
\end{itemize}

\subsection{Invariantes}

\begin{itemize}
  \item La API externa no redefine reglas del dominio.
  \item Todo resultado vigente usado por scoring debe tener fuente identificable.
  \item El frontend no consume directamente la estructura del proveedor externo.
  \item Un partido no puede tener simultaneamente mas de un resultado vigente a efectos del sistema.
\end{itemize}

\subsection{Preguntas abiertas}

\begin{itemize}
  \item Validar un match \texttt{FINISHED} real de World Cup 2026 cuando existan los primeros partidos finalizados (smoke test).
  \item Confirmar si el sistema recibira tambien eventos de fixture y cambios de kickoff desde la misma integracion o desde una fuente separada.
\end{itemize}

\subsection{Decisiones pendientes}

\begin{itemize}
  \item \texttt{Decision pendiente} \texttt{No bloqueante}: validar limites reales de plan (scores delayed / schedules delayed) y su impacto operativo durante el torneo.
\end{itemize}
